\documentclass[11pt,a4paper]{article}
\usepackage[margin=22mm]{geometry}
\usepackage{amsmath,booktabs,multicol,tikz,fancyhdr}
\pagestyle{fancy}
\fancyhf{}
\fancyhead[L]{Modul Matematika}
\fancyhead[R]{Kelas X}
\fancyfoot[C]{\thepage}
\newcommand{\tujuan}[1]{\noindent\fbox{\parbox{0.94\linewidth}{#1}}}
\begin{document}
\begin{center}{\LARGE Persamaan Kuadrat}\end{center}
\tujuan{Tujuan: memahami bentuk, akar, dan grafik persamaan kuadrat.}
\section{Materi dan Rumus}
\begin{multicols}{2}
Persamaan kuadrat memiliki bentuk $ax^2+bx+c=0$.
\[x=\frac{-b\pm\sqrt{b^2-4ac}}{2a}\]
\columnbreak
\begin{tabular}{lr}\toprule
Koefisien & Nilai\\\midrule
$a$ & 1\\ $b$ & -3\\ $c$ & 2\\\bottomrule
\end{tabular}
\end{multicols}
\section{Grafik}
\begin{center}\begin{tikzpicture}[scale=0.7]
\draw[->] (-1,0)--(4,0) node[right] {$x$};
\draw[->] (0,-2)--(0,4) node[above] {$y$};
\draw[blue,thick,domain=-0.3:3.3,samples=60] plot (\x,{\x*\x-3*\x+2});
\end{tikzpicture}\end{center}
\newpage
\section{Latihan}
Halaman ini harus tetap menjadi halaman kedua.
\begin{enumerate}
\item Tentukan akar $x^2-3x+2=0$.
\item Jelaskan hubungan diskriminan dengan banyaknya akar real.
\end{enumerate}
\end{document}
